\section{Conventions and initial reductions}\label{sec:preliminaries}

We fix the coefficient and automorphic conventions before
introducing deformation spaces. We then remove the virtually
solvable case and specify the trace loci used throughout the
proof. Table~\ref{tab:conventions} collects the normalizations that
link deformation theory, completed forms and classicality.

\begin{table}[htbp]
\centering\small
\renewcommand{\arraystretch}{1.2}
\begin{tabular}{@{}p{0.23\textwidth}p{0.70\textwidth}@{}}
\toprule
Object & Convention \\
\midrule
Frobenius & Geometric by default; local reciprocity sends a uniformizer
to geometric Frobenius. The tame relation with exponent \(q_v\)
uses arithmetic Frobenius explicitly. \\
Weights & \(\eps=x|x|\), cyclotomic exponent \(1\),
\(\HT(\eps)=\{-1\}\), and \(|2|=2^{-1}\). \\
Determinants and blocks & \(\chi=\det r\), completed-form central
character \(\psi=\chi\eps\), local block parameter \(r\otimes\eps\)
with determinant \(\chi\eps^2\). \\
Ordinary target & At a locally reducible de Rham target choose the line character
\(\eta_1\) with exponent \(a>b\), where
\(\eta_2=\chi/\eta_1\) has exponent \(b\). The favorable upper-Borel,
unnormalized inducing pair is \((\eta_2\eps,\eta_1)\). \\
Algebraic factor & \(\operatorname{Sym}^{a-b-1}\det^{b+1}\).
The continuous family retaining \(\eta_1\) need not be de Rham. \\
Curve coefficients & \(K=k((t))\) has its Laurent topology;
\(\mathcal L=\widehat{\OO[[z]]_{(\varpi)}}\), with \(z\mapsto t\).
The auxiliary ultrapower in Section~\ref{sec:propagation} is specified
separately. \\
Dual actions & Group duals use the indicated contragredient action.
On scalar eigenspaces the commuting Hecke, torus and \(U\) operators
are transposed without inversion, as in Section~\ref{sec:classicality}. \\
Support & Images in pseudodeformation spectra are closed images;
pro-modular support containments are set-theoretic. \\
\bottomrule
\end{tabular}
\caption{Conventions used at the interfaces of the proof. The
character orientation in the table concerns the chosen target and
the normalized dense points, not every member of the family.}
\label{tab:conventions}
\end{table}

\subsection{Coefficients and cyclotomic exponents}

Write \(\OO\) for the ring of integers in a sufficiently large
finite extension \(E/\mathbb Q_2\), \(\varpi\) for a uniformizer,
and \(k\) for its residue field. Finite coefficient extensions
are allowed. A lattice means a finite free \(\OO\)-module of
full rank. The cyclotomic character is denoted by \(\eps\).
We count weights as \emph{cyclotomic exponents}, so that
\(\eps\) has exponent \(1\). In the Hodge--Tate convention
\(\HT(\eps)=\{-1\}\), exponents are the negatives of
Hodge--Tate weights.

We use geometric Frobenius and local reciprocity sending a
uniformizer to geometric Frobenius. Thus, as a character of
\(\mathbb Q_2^\times\),
\[
 \eps=x|x|,\qquad x(a)=a.
\]
The determinant of a Galois representation is denoted by
\(\chi\). For a representation in Theorem~\ref{thm:main},
de Rham class field theory gives
\[
 \chi=\chi_{\mathrm{fin}}\eps^w
\]
for a finite-order character \(\chi_{\mathrm{fin}}\) and an
integer \(w\). Indeed, Tate's local algebraicity theorem
\cite[Chapter~III, Section~1.2]{Serre1989Abelian} gives the integral
cyclotomic exponent at~\(2\). Removing this power leaves a character
locally constant at every finite place. Finite ramification and global
reciprocity make it a character of a finite ray class group, hence of
finite order.
In the convention for quaternionic forms
used below, the central character is
\[
 \psi=\chi\eps.
\]
The normalized local block parameter attached to a Galois
representation \(\rho\) is that of \(\rho\otimes\eps\).
These conventions agree with \cite[Sections~3.1--3.5]{Pan2022}.

For the statement of Theorem~\ref{thm:main}, one may use the arithmetic
normalization \(V_f\) of the representation of a primitive eigenform
of weight \(k\) and nebentype \(\psi_f\): at a good prime \(q\),
\[
 \tr V_f(\operatorname{Frob}^{\mathrm{ar}}_q)=a_q(f),\qquad
 \det V_f=\psi_f\eps^{k-1}.
\]
These trace and determinant formulas use the normalization of
\cite[Theorem~6.1]{DeligneSerre1974}; arithmetic Frobenius is fixed
there in Section~3.1. The cyclotomic exponents of \(V_f\) are
\(\{0,k-1\}\). Evaluating the same
representation at geometric Frobenius inverts its eigenvalues; it
does not change the representation into its dual. If a construction
instead produces \(V_f^\vee\), the wedge pairing gives
\[
 V_f^\vee\simeq V_f\otimes\psi_f^{-1}\eps^{1-k}
       \simeq V_g\otimes\eps^{1-k},
\]
where \(g\) is the primitive form associated with the finite-character
twist \(f\otimes\psi_f^{-1}\). Thus an expression
\(r\simeq V_f^\vee\eps^n\) becomes
\(r\simeq V_g\eps^{n+1-k}\). The up-to-Tate-twist conclusion of the
main theorem is unchanged by this conversion.

\begin{lemma}\label{lem:finite-coefficients}
A continuous finite-dimensional representation of a profinite
group on \(\overline{\mathbb Q}_2^n\) is defined over a finite
extension of \(\mathbb Q_2\) and admits a stable lattice there.
\end{lemma}
\begin{proof}
There are countably many finite subextensions of
\(\overline{\mathbb Q}_2/\mathbb Q_2\). For each one \(E\),
the set of group elements whose representation matrices
lie in \(\GL_n(E)\) is a closed subgroup. These subgroups
cover the profinite group. By Baire's theorem, one has
nonempty interior, hence is open. Enlarge \(E\) to contain
the finitely many matrix entries of coset representatives.
Then every matrix lies in that finite extension.

The image is compact. Its entries, and the entries of its
inverses, have bounded valuations. The \(\OO\)-submodule
generated by the orbit of \(\OO^n\) therefore lies between
two scalar multiples of \(\OO^n\). It is a stable lattice.
\end{proof}

We apply Lemma~\ref{lem:finite-coefficients} whenever a
characteristic-zero representation is first introduced.
Over a curve field, the analogous compact-image lattice
argument gives an integral realization over \(k'[[t]]\).

\subsection{The virtually solvable case}

For a representation over any field, after extending scalars
to an algebraic closure, call it \emph{virtually solvable}
if the identity component of the reduced Zariski closure of its image
is solvable. The closure here is taken over the algebraic closure,
so that the resulting algebraic group is smooth. This property is unchanged by restriction to
a finite-index subgroup.

\begin{lemma}\label{lem:solvable-reduction}
Theorem~\ref{thm:main} holds for virtually solvable
representations. A representation satisfying its hypotheses
and not virtually solvable remains absolutely irreducible
and not virtually solvable over every finite extension of
\(\mathbb Q\).
\end{lemma}
\begin{proof}
Let \(H\) be the algebraic monodromy group of an irreducible
two-dimensional representation over an algebraic closure.
Lie--Kolchin and the structure of connected solvable groups
\cite[Theorems~17.4 and~16.33]{MilneAlgebraicGroups} give the following
alternatives. If \(H^\circ\) is solvable and has a nontrivial unipotent
radical, the latter has a unique fixed line; normality
makes that line \(H\)-stable, contradicting irreducibility.
The remaining possibilities are that \(H^\circ\) is scalar,
or that its two distinct character lines are permuted by
\(H\).

In the scalar case the projective image is finite. After a
finite extension of the base field, the representation is
a scalar character on a two-dimensional space. Its two
Hodge--Tate weights are equal, which is excluded. In the
second case the action on the two lines gives a quadratic
extension \(L/\mathbb Q\) and
\[
 r\simeq\Ind_{G_L}^{G_{\mathbb Q}}\theta
\]
for a continuous character \(\theta\).
The character is de Rham and finitely ramified, hence is the
\(2\)-adic avatar of an algebraic Hecke character, by the
following rank-one consequence of Tate's theorem and the construction
in \cite[Chapter~III, Sections~1.2 and~2.3]{Serre1989Abelian}.
The integral algebraic exponents on open local units obey the global
unit relation. Principal ideals in a suitable ray subgroup therefore
receive the prescribed algebraic monomials. A positive power of each
of finitely many ray-class representatives is principal in that
subgroup, so its character value is algebraic. Adjoining these values
places all ideal-character values in one number field and constructs
the algebraic Hecke character. This argument does not assume that the
given character was initially rational over a number field.
If \(L\)
is real quadratic, its algebraic exponents at the two real
embeddings are equal by the unit theorem: a unit of infinite
order forces their difference to vanish in the algebraic
character relation. This again contradicts regularity.
For imaginary quadratic \(L\), the theta-series construction
gives the required cuspidal eigenform, with its algebraic
Tate normalization; equivalently, use quadratic automorphic
induction \cite[Proposition~12.1]{JacquetLanglands1970}.

Finally, restriction to finite index preserves \(H^\circ\).
If a restriction of a non-virtually-solvable representation
were reducible, its identity component would lie in a
Borel subgroup and be solvable. This is impossible.
\end{proof}

Henceforth the target \(r\) is not virtually solvable.
In particular, its restriction to every field used below
is absolutely irreducible.

\subsection{Fields, deformation spaces, and trace loci}

All auxiliary base fields \(F\) are totally real, solvable
over \(\mathbb Q\), of even degree, and split completely at
\(2\). We may replace them by Galois closures. Let
\(S\) contain the real places, all places above \(2\), and
a finite set of other places. Set
\[
 d=[F:\mathbb Q],\qquad
 u=\#\{v\in S:v\nmid 2,\ v\text{ finite}\},\qquad
 s=2d+u.
\]
The group \(G_{F,S}\) is the Galois group of the maximal
extension of \(F\) unramified outside \(S\).

A two-dimensional pseudorepresentation means a determinant
in the sense of Chenevier, with its determinant character
retained. In particular, no division by \(2\) is used to
recover a determinant from traces. Fixing \(\chi\) and a
semisimple residual representation gives a complete local
pseudodeformation ring, denoted by \(R_F^{\mathrm{ps},\chi}\).
Framed representation rings and schemes, including those
with specified local invariant lines, will be constructed
in Section~\ref{sec:deformation}. Whenever a component is
mapped to a pseudodeformation spectrum, its \emph{image}
means the closure of its image unless stated otherwise.

There is a unique quaternion algebra \(D_F/F\), up to
isomorphism, ramified at all real places and at no finite
place. The evenness of \(d\) makes this possible, by the classification by
ramification sets \cite[Main Theorem~14.6.1]{Voight2021}.
A completed Hecke algebra of definite \(D_F^\times\)-forms,
with central character \(\psi=\chi\eps\), defines a closed
trace locus in \(\Spec R_F^{\mathrm{ps},\chi}\).
We call such a locus \emph{pro-modular}. The tame level is
open, is hyperspecial away from \(S\), and is unrestricted
at the places above \(2\). A closed trace locus is
\emph{potentially pro-modular} if its restriction becomes
pro-modular after a further solvable totally real extension
split completely at \(2\). All support containments in this
terminology are set-theoretic.

A \emph{characteristic-two curve} is a one-dimensional
domain quotient of a complete local \(k\)-algebra. After
finite extension, its normalization has the form
\(k'[[t]]\); its generic point has residue field
\(k'((t))\). We also use the word curve for this generic
point and its associated representation when there is
no ambiguity. Whenever such a representation is used,
we choose a compact-image lattice. Agreement of traces
and determinants modulo the maximal ideal determines
the semisimple residual representation, but need not
determine a residual extension class. The latter
distinction will be important in the marked-line
connectedness argument.

The reductions now leave a fixed determinant and a target that remains
absolutely irreducible after every finite restriction. We next construct
the spaces in which its trace, and when needed its local line character,
can vary. The distinction between these data is essential: a trace
space is sufficient for completed-form occurrence, whereas the final
reducible classicality argument also needs the chosen character.
